\documentclass[12pt,a4paper]{article}
\usepackage[utf8]{vietnam}
\usepackage{amsmath,amssymb}
\usepackage{geometry}
\geometry{
  top=12mm,
  bottom=12mm,
  left=14mm,
  right=14mm,
  footskip=6mm,
  headheight=14pt
}

\usepackage{tikz}
\usetikzlibrary{arrows.meta,calc,shapes.geometric}
\usepackage[most]{tcolorbox}
\usepackage{fontawesome}
\usepackage{fancyhdr}
\usepackage{enumitem}

\pagestyle{fancy}
\fancyhf{}
\fancyhead[L]{\footnotesize\color{teal!80!black}\faBook\ \textbf{SỔ TAY KIẾN THỨC TOÁN 3}}
\fancyhead[R]{\footnotesize\color{teal!80!black}\textbf{Lớp Toán Thầy Lĩnh -- loptoanthaylinh.vn}}
\fancyfoot[C]{\footnotesize\color{gray}-- Sổ tay kiến thức: Bài 26. Chia số có hai chữ số cho số có một chữ số --}
\renewcommand{\headrulewidth}{0.4pt}
\renewcommand{\footrulewidth}{0pt}

\setlength{\parindent}{0pt}
\setlength{\parskip}{0pt}

\begin{document}

% --- BANNER BÀI HỌC ---
\begin{tcolorbox}[
  enhanced,
  colback=teal!8!white,
  colframe=teal!70!black,
  arc=3mm, boxrule=1.2pt,
  left=3mm, right=3mm, top=2mm, bottom=2mm,
  halign=center
]
  {\LARGE\bfseries\color{teal!90!black}\faBookmark\ BÀI 26: CHIA SỐ CÓ HAI CHỮ SỐ CHO SỐ CÓ MỘT CHỮ SỐ}\\[1mm]
  {\itshape\footnotesize\color{black!75} Kĩ thuật đặt tính chia cột, tiến trình 3 nhịp Chia - Nhân - Trừ từ trái sang phải trên 1 trang A4}
\end{tcolorbox}

\vspace{-2mm}

\noindent
\begin{minipage}[t]{0.49\textwidth}\vspace{0pt}
\begin{tcolorbox}[
  enhanced,
  colback=blue!4!white,
  colframe=blue!75!black,
  arc=2mm, boxrule=0.8pt,
  left=2mm, right=2mm, top=1.8mm, bottom=1.8mm,
  title={\footnotesize\bfseries\color{white}\faShareAlt\ 1. Quy trình Đặt tính rồi tính}
]
{\footnotesize
\begin{itemize}\setlength{\itemsep}{1pt}\setlength{\leftskip}{-2.5mm}
  \item \textbf{Thứ tự thực hiện:} Khác với $(+, -, \times)$, phép chia thực hiện lần lượt từ \textbf{trái sang phải} (từ hàng chục đến hàng đơn vị).
  \item \textbf{Nhịp 3 bước ở mỗi lượt chia:}
  \begin{enumerate}\setlength{\itemsep}{0.5pt}\setlength{\leftskip}{-2.5mm}
    \item \textbf{CHIA:} Lấy chữ số chia cho số chia để tìm chữ số của thương.
    \item \textbf{NHÂN:} Lấy thương vừa tìm được nhân ngược lại với số chia.
    \item \textbf{TRỪ:} Lấy số bị chia trừ đi tích để tìm số dư của lượt đó.
  \end{enumerate}
\end{itemize}
}
\end{tcolorbox}
\end{minipage}
\hfill
\begin{minipage}[t]{0.49\textwidth}\vspace{0pt}
\begin{tcolorbox}[
  enhanced,
  colback=teal!4!white,
  colframe=teal!75!black,
  arc=2mm, boxrule=0.8pt,
  left=2mm, right=2mm, top=1.8mm, bottom=1.8mm,
  title={\footnotesize\bfseries\color{white}\faSliders\ 2. Kĩ thuật hạ chữ số hàng đơn vị}
]
{\footnotesize
\begin{itemize}\setlength{\itemsep}{1pt}\setlength{\leftskip}{-2.5mm}
  \item Sau khi chia xong hàng chục (dư $0$ hoặc dư $r$):
  \begin{itemize}\setlength{\itemsep}{0.5pt}\setlength{\leftskip}{-2.5mm}
    \item \textbf{Hạ} chữ số hàng đơn vị xuống bên phải số dư để tạo thành số mới.
    \item Thực hiện tiếp nhịp 3 bước (Chia $\to$ Nhân $\to$ Trừ) cho lượt thứ hai.
  \end{itemize}
  \item Số dư cuối cùng của lượt hai chính là \textbf{số dư của toàn bộ phép chia}.
\end{itemize}
}
\end{tcolorbox}
\end{minipage}

\vspace{-2.5mm}
% --- KHỐI 3: SƠ ĐỒ TƯ DUY TIKZ ---
\noindent
\begin{tcolorbox}[
  enhanced,
  colback=orange!4!white,
  colframe=orange!80!black,
  arc=2mm, boxrule=0.8pt,
  left=2mm, right=2mm, top=1.5mm, bottom=1.5mm,
  title={\footnotesize\bfseries\color{white}\faSitemap\ SƠ ĐỒ TƯ DUY: TIẾN TRÌNH 2 LƯỢT CHIA CỘT TỪ TRÁI SANG PHẢI}
]
\centering
\begin{tikzpicture}[
  box/.style={draw, rounded corners=2pt, align=center, font=\scriptsize\bfseries, inner sep=3pt, minimum height=0.85cm},
  arr/.style={-{Stealth[length=2.5mm]}, thick}
]
  \node[box, fill=blue!12, draw=blue!80!black, text width=3.3cm] (d1) at (0, 0) {LƯỢT 1: HÀNG CHỤC\\Chia chục $\to$ Nhân $\to$ Trừ\\Tìm số dư hàng chục};
  \node[box, fill=teal!15, draw=teal!80!black, text width=3.3cm] (d2) at (4.7, 0) {HẠ ĐƠN VỊ XUỐNG\\Ghép với số dư lượt 1\\Tạo số bị chia lượt 2};
  \node[box, fill=purple!12, draw=purple!80!black, text width=3.3cm] (d3) at (9.4, 0) {LƯỢT 2: HÀNG ĐƠN VỊ\\Chia đơn vị $\to$ Nhân $\to$ Trừ\\Ghi thương \& Số dư cuối};

  \draw[arr, blue!80!black] (d1) -- (d2);
  \draw[arr, teal!80!black] (d2) -- (d3);
\end{tikzpicture}
\end{tcolorbox}

\vspace{-2.5mm}
% --- KHỐI 4: VÍ DỤ & MẸO TRÁNH BẪY ---
\noindent
\begin{minipage}[t]{0.49\textwidth}\vspace{0pt}
\begin{tcolorbox}[
  enhanced,
  colback=green!4!white,
  colframe=green!60!black,
  arc=2mm, boxrule=0.8pt,
  left=2mm, right=2mm, top=1.8mm, bottom=1.8mm,
  title={\footnotesize\bfseries\color{white}\faCheckCircle\ Ví dụ thực hành chuẩn mực}
]
{\footnotesize
\textbf{1. Chia hết ($48 : 2$):}\\
$\bullet\ 4$ chia $2$ được $2$, viết $2$; $2 \times 2 = 4$; $4 - 4 = 0$.\\
$\bullet$ Hạ $8$; $8$ chia $2$ được $4$, viết $4$; $4 \times 2 = 8$; $8 - 8 = 0$.\\
$\implies 48 : 2 = \mathbf{24}$.\\[1mm]
\textbf{2. Chia có dư ($75 : 3$ và $58 : 5$):}\\
$\bullet\ 75 : 3 = \mathbf{25}$ (chia hết).\\
$\bullet\ 58 : 5$: $5 : 5 = 1$; hạ $8$, $8 : 5 = 1$ dư $3 \implies \mathbf{11\text{ (dư 3)}}$.
}
\end{tcolorbox}
\end{minipage}
\hfill
\begin{minipage}[t]{0.49\textwidth}\vspace{0pt}
\begin{tcolorbox}[
  enhanced,
  colback=red!4!white,
  colframe=red!70!black,
  arc=2mm, boxrule=0.8pt,
  left=2mm, right=2mm, top=1.8mm, bottom=1.8mm,
  title={\footnotesize\bfseries\color{white}\faExclamationTriangle\ Mẹo tránh bẫy trắc nghiệm}
]
{\footnotesize
\begin{itemize}\setlength{\itemsep}{1.5pt}\setlength{\leftskip}{-2.5mm}
  \item \textbf{Bẫy chia từ hàng đơn vị:}\\
  Cộng, trừ, nhân đều làm từ phải sang trái, nhưng phép chia BẮT BUỘC phải làm từ \textbf{trái sang phải} (hàng chục trước).
  \item \textbf{Bẫy quên hạ chữ số hàng đơn vị:}\\
  Sau khi chia xong hàng chục, phải hạ hàng đơn vị xuống chia tiếp lượt hai, không được kết luận luôn thương chỉ có 1 chữ số!
\end{itemize}
}
\end{tcolorbox}
\end{minipage}


\end{document}
